% Append this section before \end{document} in the master AION document.
\section{Signed Model Catalogue and Evaluation Fabric}
\label{sec:signed-model-catalogue}

The AION brain is deliberately independent of any individual language model.  This
creates a second requirement: a replacement model must not be trusted merely because
it can be downloaded or because a provider claims that it is superior.  Pilot therefore
uses a customer-controlled model catalogue that establishes identity, compatibility,
licensing, measured quality and safe promotion before a model can serve live work.

The catalogue is implemented by the \texttt{SignedModelCatalogue} authority.  Its state
is separate from the canonical Business Map and from customer operational records.
Installing, replacing, blocking or removing a model consequently cannot erase or rewrite
what AION has learned about the business.  Models remain replaceable computational tools;
the brain retains identity, policy, memory, evidence, permissions and outcome history.

\subsection{Signed and Immutable Model Identity}

Each model release is described by a versioned manifest containing:
\begin{itemize}
  \item a stable model identifier and immutable version;
  \item the provider and inference format;
  \item the artefact's SHA--256 digest and exact byte length;
  \item the licence identifier and explicit commercial-use and redistribution flags;
  \item approved download mirrors;
  \item supported processor architectures, accelerator classes and inference engines;
  \item minimum memory and storage requirements; and
  \item the issuing public key, manifest hash and Ed25519 signature.
\end{itemize}

The manifest hash is calculated over canonical JSON before the signature is applied.
Registration succeeds only when the issuer is present in the local trust set, the hash
matches and the signature verifies.  A second manifest cannot silently reuse the same
model identifier and version with different content.  This prevents a mutable download
location or provider catalogue entry from redefining an already approved model release.

Revocation is also signed by the original issuer and must include a human-readable
reason.  A revoked model is marked blocked.  If it is the active champion and a previous
qualified champion exists, the route is rolled back automatically.

\subsection{Installation and Hardware Qualification}

Before installation, Pilot compares the manifest with the coarse hardware profile
created during sovereign-brain setup.  The preflight checks available memory and storage,
processor architecture, accelerator class and selected inference engine.  It then streams
the candidate artefact through SHA--256 verification and checks its byte length, avoiding
the need to load a large weight file into memory.  Commercial deployment is rejected when
the manifest does not explicitly permit commercial use.

The resulting installation decision is fail-closed and reports concrete reasons such as
insufficient memory, an unsupported architecture, an incompatible inference engine, a
licence restriction or an artefact-integrity mismatch.  Compatibility is therefore a
measured deployment property rather than an assumption based on a model name.

\subsection{Seven-Dimensional Evaluation}

A model cannot be qualified without a complete benchmark record.  Every benchmark suite
must report all of the following dimensions:
\begin{enumerate}
  \item task performance;
  \item domain performance;
  \item safety performance;
  \item governed tool-use performance;
  \item latency;
  \item monetary cost; and
  \item energy use.
\end{enumerate}

Quality, safety and tool-use dimensions use minimum acceptance thresholds.  Latency,
cost and energy use maximum thresholds.  The benchmark includes a suite identifier,
provenance and a content hash so that a result can be traced to its evaluation method.
A model is labelled \emph{qualified} only when it is compatible with the intended machine
and engine and the latest public benchmark meets every customer threshold.  Otherwise it
is labelled \emph{blocked}; unevaluated valid releases remain \emph{experimental}.

Public repeatable benchmark results and customer-private outcome scores are intentionally
separated.  Public results may appear in the administrator catalogue.  Private results
are stored in a separate owner-readable file, require a customer identifier and never
appear in the administrator projection.  This permits AION to learn which model works
best for a particular business without leaking its performance, workloads or outcomes
into a shared model registry.

\subsection{Champion, Challenger and Bounded Promotion}

Each task class may have an active champion and one challenger.  A canary can direct no
more than 20 percent of eligible work to the challenger and may collect no more than
1,000 paired requests.  Both models must already be qualified, and a caller cannot replace
the recorded champion merely by naming a different one when opening a canary.

For each sample, the catalogue records comparable champion and challenger measurements.
At the request ceiling it calculates the mean for every configured threshold.  The
challenger must both meet the absolute customer threshold and equal or outperform the
champion: higher is better for quality dimensions and lower is better for latency, cost
and energy.  Failure on any dimension rejects the challenger.  Success promotes it while
retaining the previous champion as a rollback target.

This mechanism deliberately avoids ``majority confidence'' or model popularity as a
promotion rule.  A candidate becomes the default only after bounded, task-specific,
customer-defined evidence demonstrates that it is not inferior on any required measure.

\subsection{Regression, Drift and Provider Change}

After promotion, Pilot can compare current measurements with a stored baseline and a
per-metric tolerance.  A quality reduction, a latency/cost/energy increase, or a change
to the provider-behaviour fingerprint constitutes drift.  Detected drift blocks the model
and triggers rollback where a previous champion is available.  The report preserves the
affected metric names and whether provider behaviour changed, without storing customer
prompt or result content.

\subsection{Licence-Controlled Mirroring}

Offline and private deployments do not mirror an unrestricted catalogue.  Pilot constructs
a bounded mirror plan from the customer's active deployment profiles.  It includes only
registered, non-blocked models whose licence explicitly permits redistribution, and only
the artefact identity and approved mirror required by those profiles.  Inactive profiles,
revoked releases and non-redistributable artefacts are excluded.  This reduces supply-chain
exposure, disk use and licence risk while retaining the packages required for recovery.

\subsection{Administrator Experience}

The administrator view groups releases into qualified, blocked and experimental lists.
For each release it shows provider, licence, qualification record, public benchmarks and
revocation state, together with active task routes and canary status.  It explicitly states
that customer-private outcomes are not exposed.  Operational controls can therefore answer
four distinct questions: \emph{What is available?  What can run here?  What has passed our
tests?  What is presently serving this task?}

\subsection{Security and Failure Properties}

The implementation enforces the following invariants:
\begin{itemize}
  \item an untrusted or invalidly signed manifest cannot enter the catalogue;
  \item a version cannot be redefined after registration;
  \item an artefact cannot pass installation when its bytes differ from the manifest;
  \item incomplete benchmark suites cannot qualify a model;
  \item customer-private outcomes cannot enter the public administrator projection;
  \item an inferior challenger cannot become the default;
  \item canary traffic and duration are bounded;
  \item revocation or measured drift blocks continued promotion and supports rollback;
  \item mirroring requires both active demand and explicit redistribution permission; and
  \item no catalogue operation modifies canonical business state.
\end{itemize}

\subsection{Implementation Status}

This phase is implemented in AION Sovereign Brain Bootstrap~0.7.0 and Pilot
Fabric~0.59.0.  Automated verification covers trusted signing, tamper rejection,
immutability, hardware and artefact checks, benchmark completeness, public/private
separation, qualification, bounded promotion, inferior-candidate rejection, manual and
automatic rollback, provider drift, signed revocation, licence-aware mirroring and
business-state independence.

Production distribution still requires organisation-specific issuer governance, secure
mirror operation, platform package signing and real hardware/model qualification.  Those
are deployment gates rather than permission to bypass the catalogue.  A model without
valid identity, compatibility and measured acceptance remains unavailable for live work.
